% Preámbulo: tipografía, paleta monocromática y macros de notación.
\usepackage[T1]{fontenc}
\usepackage[utf8]{inputenc}
\usepackage{lmodern}
\usepackage{amsmath}
\usepackage[spanish,es-noquoting,es-noshorthands,es-tabla]{babel}
\usepackage[a4paper,margin=2.6cm,headsep=14pt]{geometry}
\usepackage{microtype}
\usepackage{booktabs}
\usepackage{longtable}
\usepackage{array}
\usepackage{xcolor}
\usepackage{fancyhdr}
\usepackage{titlesec}
\usepackage{enumitem}
\usepackage{fancyvrb}
\usepackage{tikz}
\usetikzlibrary{automata,positioning,arrows.meta,calc}
\usepackage[hidelinks,bookmarks=true]{hyperref}

% --- Paleta monocromática: sólo grises, con un único acento. -----------------
\definecolor{tinta}{gray}{0.13}
\definecolor{grisMedio}{gray}{0.42}
\definecolor{grisSuave}{gray}{0.70}
\definecolor{grisFondo}{gray}{0.955}
\definecolor{acento}{HTML}{2F4858}

\color{tinta}

% --- Títulos ----------------------------------------------------------------
\titleformat{\section}{\Large\bfseries\color{acento}}{\thesection}{0.7em}{}
\titleformat{\subsection}{\large\bfseries\color{tinta}}{\thesubsection}{0.7em}{}
\titleformat{\subsubsection}{\bfseries\color{grisMedio}}{\thesubsubsection}{0.7em}{}
\titlespacing*{\section}{0pt}{2.2ex plus 1ex minus .2ex}{1.2ex plus .2ex}

% --- Encabezado y pie -------------------------------------------------------
\setlength{\headheight}{14pt}
\pagestyle{fancy}
\fancyhf{}
\renewcommand{\headrulewidth}{0.4pt}
\renewcommand{\headrule}{\hbox to\headwidth{\color{grisSuave}\leaders\hrule height \headrulewidth\hfill}}
\fancyhead[L]{\footnotesize\color{grisMedio}Compiladores e Intérpretes — Etapa 1}
\fancyhead[R]{\footnotesize\color{grisMedio}\nouppercase{\leftmark}}
\fancyfoot[C]{\footnotesize\color{grisMedio}\thepage}

% --- Bloques de código ------------------------------------------------------
\DefineVerbatimEnvironment{codigo}{Verbatim}%
  {fontsize=\small,xleftmargin=1.1em,commandchars=\\\{\}}
\DefineVerbatimEnvironment{codigoLlano}{Verbatim}%
  {fontsize=\small,xleftmargin=1.1em}

% --- Macros de notación -----------------------------------------------------
\newcommand{\bs}{\textbackslash}
\newcommand{\caret}{\textasciicircum}
\newcommand{\lex}[1]{\texttt{#1}}
\newcommand{\tok}[1]{\textnormal{\textsf{\small #1}}}
% --- Cardinalidad como exponente propiamente dicho -------------------------
% \sups{c}    exponente suelto:  ( ... )^*
% \ere{X}{c}  símbolo con su exponente:  Digito^{1-9}
\newcommand{\sups}[1]{\textsuperscript{\normalfont\texttt{\footnotesize #1}}}
\newcommand{\ere}[2]{\texttt{#1}\sups{#2}}
\newcommand{\rep}[2]{\sups{#1--#2}}
\newcommand{\alt}{\ensuremath{\;|\;}}

% Nota al pie de notación, repetible: se recuerda la convención donde aparece.
\newcommand{\notaKleene}{\footnote{\textbf{Notación.} \ere{X}{*} = cero o más
repeticiones de \lex{X} (estrella de Kleene); \ere{X}{+} = una o más;
\ere{X}{?} = cero o una; \ere{X}{1--9} = entre 1 y 9 repeticiones, \emph{ambos
extremos inclusive}; \lex{A | B} = alternativa. La tabla completa está en la
sección~\ref{sec:notacion}.}}

\newenvironment{er}{%
  \par\smallskip\begingroup\ttfamily\small\parindent=0pt
  \leftskip=1.2em\obeylines}{\par\endgroup\smallskip}

\setlist[itemize]{leftmargin=1.3em,itemsep=0.25ex,topsep=0.5ex}
\setlist[description]{leftmargin=1.3em,itemsep=0.35ex,topsep=0.6ex}
\renewcommand{\arraystretch}{1.25}
